\section{Related Work}

\paragraph{Synthetic cold-start evaluation in DTI.}
Machine-learning models for DTI and affinity prediction are commonly evaluated on BindingDB- and DAVIS-derived random, cold-drug, cold-target, or blind splits \citep{liu2007bindingdb,davis2011kinase,ahmed2024dtilm,luo2024druglamp,ouyang2024pmmr,liu2025spdti,yu2025gsdti,zhao2025coldstartcpi}. These protocols are useful for controlled comparison, but they eliminate different forms of chemical, protein, and pair overlap and therefore do not represent a single validation setting. For this reason, results obtained on one synthetic split should not automatically be interpreted as evidence of general robustness.

\paragraph{Validation design in chemoinformatics.}
More broadly in chemoinformatics, validation design and test-set construction are known to influence apparent model performance. Reanalysis of large bioactivity benchmarks, increasingly challenging potency-prediction partitions, and transfer-learning studies for binary drug--target interactions have all shown that the relative standing of deep and classical models can change materially when the evaluation setting changes \citep{robinson2020validating,rodriguezperez2021evaluation,tanoori2021binding}. Our study follows this validation-oriented tradition by asking how split semantics alter conclusions in a DTI setting built around widely used BindingDB and DAVIS resources.

\paragraph{Benchmark and resource studies under distribution shift.}
Recent benchmark studies in molecular machine learning and related biomedical prediction problems have shown that scaffold, temporal, and domain-shifted evaluation can change model ranking and alter the practical interpretation of benchmark results \citep{ektefaie2024generalizability,shen2025ddiben,delafuente2025inductive}. This literature has shifted attention from purely i.i.d.\ comparison toward reusable benchmark resources that expose how conclusions depend on data partitioning. Our work extends that perspective to DTI by coupling multiple benchmark families with exported split manifests and a standalone evaluator.

\paragraph{Positioning.}
Prior DTI papers typically optimize performance within one benchmark design. In contrast, our contribution is a reproducible case study of ranking stability across synthetic cold-start and temporal/provenance benchmarks, together with overlap-sensitive robustness analyses and reusable evaluation resources. The emphasis is therefore on validation methodology and benchmark semantics rather than on proposing a new predictive architecture.
